\ifCLASSOPTIONcompsoc
\IEEEraisesectionheading{\section{Introduction}\label{sec:introduction}}
\else
\section{Introduction}
\label{sec:introduction}
\fi

\IEEEPARstart{T}{his} section is currently an outline. Each bullet is meant
to become one or two sentences of running text.

\subsection{Motivation}
\begin{itemize}
  \item Non-player characters (NPCs) are every character in a game that the
        player does not control: enemies, companions, shopkeepers, crowds.
  \item They largely decide whether a game world feels alive or empty; a
        guard who walks into a wall breaks immersion faster than a weak
        texture.
  \item Game AI differs from academic AI: the goal is not to win or to be
        optimal, but to create an entertaining and believable experience for
        the player~\cite{yannakakis2018,millington2009}.
  \item Laird and van Lent called interactive games the ``killer
        application'' for human-level AI because they demand real-time,
        social and adaptive behavior~\cite{laird2001}.
  \item The field is changing fast: hand-authored systems from the 1990s and
        2000s now sit next to machine learning and large language models
        (LLMs)~\cite{gallotta2024}.
\end{itemize}

\subsection{What Makes a Good NPC?}
\begin{itemize}
  \item \emph{Believability}: the character appears to have its own goals,
        emotions and awareness, even if its intelligence is
        limited~\cite{bates1994,loyall1997}.
  \item Warpefelt and Verhagen describe believability as fitting the
        player's expectations for that kind of character in that
        world~\cite{warpefelt2017}.
  \item \emph{Fun and fairness}: an enemy should be challenging but
        beatable; perfect aim is rarely enjoyable~\cite{sweetser2005}.
  \item \emph{Designer control}: game designers need to shape and predict
        what NPCs do so that levels and stories work as intended.
  \item \emph{Efficiency}: AI shares a small slice of each frame with
        graphics, physics and audio, often only a few
        milliseconds~\cite{nareyek2004}.
  \item These goals often conflict, which is a recurring theme of this
        review (see Section~\ref{sec:tradeoffs}).
\end{itemize}

\subsection{Scope and Contributions}
\begin{itemize}
  \item Focus on \emph{decision making} (what an NPC does) and \emph{spatial
        reasoning} (where and how it moves); animation and graphics are out
        of scope.
  \item Covers commercial techniques documented by industry developers as
        well as academic research.
  \item Contributions: (1) an accessible overview of the main NPC AI
        techniques, (2) a side-by-side comparison along shared criteria, and
        (3) a discussion of current experiments with learning and LLM-based
        NPCs.
\end{itemize}

\subsection{Methodology}
\begin{itemize}
  \item Literature review of textbooks~\cite{millington2009,yannakakis2018},
        peer-reviewed papers (e.g.\ IEEE Transactions on Games, AIIDE) and
        industry talks from the Game Developers Conference (GDC).
  \item Industry talks are included because many commercial techniques are
        only documented there; this is noted as a limitation.
  \item Selection criteria: influence on practice (used in shipped games)
        or on research (widely cited).
\end{itemize}

\subsection{Structure of the Paper}
\begin{itemize}
  \item Section~\ref{sec:foundational}: early and foundational systems.
  \item Section~\ref{sec:modern}: modern decision-making architectures.
  \item Section~\ref{sec:perception}: perception and spatial reasoning.
  \item Section~\ref{sec:tradeoffs}: comparison, constraints and
        trade-offs.
  \item Section~\ref{sec:experiments}: current experiments with machine
        learning and LLMs.
  \item Section~\ref{sec:conclusion}: conclusion and outlook.
\end{itemize}
